\begin{helper}
Fix a history cell \(C\), a fixed-count subevent \(D\subseteq C\), a terminal
set \(U\), and one of the two colors.  Suppose that every \(\omega\in C\cap D\)
has a first-color terminal trace
\[
A(\omega)\subseteq U_{\mathrm{first}}
\]
such that, for each first-color terminal coordinate \(e\), the terminal value
of \(e\) in \(\omega\) is present exactly when \(e\in A(\omega)\).  Then the
complete first-color trace classes cover \(C\cap D\), and two distinct trace
sets define disjoint classes.
\end{helper}
\begin{proof}
Let \(U_{\mathrm{first}}\) be the first-color part of \(U\).  By hypothesis,
for each \(\omega\in C\cap D\) there is a subset
\(A(\omega)\subseteq U_{\mathrm{first}}\) recording exactly the present
first-color terminal coordinates.  Therefore \(\omega\) belongs to the branch
indexed by \(A(\omega)\), so the branches cover \(C\cap D\).

For disjointness, let \(A\) and \(B\) be two different subsets of
\(U_{\mathrm{first}}\).  Choose \(e\) in their symmetric difference.  Since
both \(A\) and \(B\) are subsets of \(U_{\mathrm{first}}\), this coordinate is
first-color and terminal.  If one sample \(\omega\) lay in both trace
branches, then the branch condition for \(A\) would say
\(X_e(\omega)=1\) exactly when \(e\in A\), while the branch condition for
\(B\) would say the same value is present exactly when \(e\in B\).  These two
truth values differ by the choice of \(e\), a contradiction.  Thus distinct
complete first-color trace branches are disjoint.
\end{proof}
